\clearpage
\onehalfspacing

\crshead{ВВЕДЕНИЕ}

Случайные числа -- это одна из основ современных криптографических
систем, алгоритмов моделирования, шифрования и даже игровых механизмов.
Качество созданных последовательностей напрямую влияет на безопасность:
так, в 2012 году из-за недостатка энтропии в Android-устройствах,
злоумышленники смогли предсказать примерно 86\% закрытых ключей
Bitcoin-кошельков, сгенерированных на уязвимых гаджетах{[}1{]}. Подобные
инциденты доказывают, что генераторы псевдослучайных чисел (ГПСЧ),
используемые в программном обеспечении, уступают по надежности
аппаратным решениям на основе физических процессов.

Физические генераторы случайных чисел (ФГСЧ) получают энтропию из
недетерминированных явлений, например, таких как тепловые шумы
резисторов, лавинные пробои полупроводников, радиоактивный распад или
даже атмосферные помехи. В отличие от математических алгоритмов,
физические методы генерации СЧ вероятно менее уязвимы к обратному
инжинирингу, но их практическая реализация сталкивается со следующими
проблемами:

\begin{enumerate}
\def\labelenumi{\arabic{enumi}.}
\item
  Более низкая скорость работы;
\item
  Потенциально высокие временные и материальные затраты на
  конструирование, установку и настройку по сравнению с программными
  ГПСЧ;
\item
  Зависимость от среды (например, шумы микрофона могут исчезнуть в
  звукоизолированном помещении);
\item
  Аппаратные искажения (смещение распределения из-за неидеальности
  компонентов).
\end{enumerate}

Цель работы -- сравнить доступные источники энтропии (аудио-, радио-,
фотонные и механические шумы) для проектирования компактного ФГСЧ с
учетом:

\begin{enumerate}
\def\labelenumi{\arabic{enumi}.}
\item
  Критериев NIST STS (тесты на частоту, последовательности, энтропию);
\item
  Стоимости и простоты реализации;
\item
  Устойчивости к внешним воздействиям.
\end{enumerate}

Актуальность моего исследования связана с растущими требованиями к
безопасности в следующих областях:

\begin{enumerate}
\def\labelenumi{\arabic{enumi}.}
\item
  Криптографические системы (генерация ключей, seed-значений);
\item
  Защищенные микропроцессоры (TPM-модули, HSM);
\item
  Специализированные устройства (криптовалютные кошельки, банковские
  терминалы).
\end{enumerate}

Проблема уязвимостей в ГСЧ остается актуальной. Атака на чипы Infineon
TPM, получившая название «ROCA», показала, что даже корпоративные
решения могут содержать фатальные недостатки.

